\documentclass[11pt,letterpaper]{article}
\usepackage{bes-whitepaper}

\begin{document}

\wptitle{Dehydration or co-ion pairing: what sets lithium selectivity\\in near-saturated brines?}
  {Separation Science}{Amanda Haes}

\hd{Research question.} In hypersaline lake brines and oilfield produced waters, lithium is a minor
ion, outnumbered by Na$^+$ and Mg$^{2+}$ by one to three orders of magnitude. Some sorbents still take
it up selectively. Why, and why does that selectivity change from one brine to the next?

\hd{Knowledge gap.} Sorbent selectivity is usually reported as a distribution coefficient measured
in a few brines, where ion activities, ion pairing and water activity all change at once, so the
mechanism rests on correlations and selectivity in a new brine cannot be predicted. Yet for the major
ions of natural brines (Na--K--Mg--Ca--Cl--SO$_4$), Pitzer models give reliable activities up to
saturation [1], and lithium-salt parameters are tabulated [4]. The solution side can therefore be
largely accounted for, and the selectivity that remains belongs to the sorbent.

\hd{Hypotheses.} \emph{H1, dehydration control:} after Pitzer activity correction, the exchange
constant $K$ for Li$^+$ against Na$^+$ or Mg$^{2+}$ depends on brine composition only through the
water activity $a_\mathrm{w}$, and the slope of $\ln K$ against $\ln a_\mathrm{w}$ counts the waters
released per exchanged ion, as expected if Mg$^{2+}$ is excluded by its far larger hydration free
energy ($-1830$ against $-475$ kJ\,mol$^{-1}$ for Li$^+$ [2]). \emph{H2, cooperative co-ion uptake:}
at fixed $a_\mathrm{w}$, selectivity also depends on which anion enters the solid with Li$^+$. H2 is
expected for a layered double hydroxide, which takes up LiCl as a pair [3], but not for a titanate
ion sieve, which exchanges Li$^+$ for H$^+$, so each sorbent is a control for the other.

\hd{Approach.} An Opentrons OT-2 prepares synthetic brines by mass in sealed vials: dilution series
of a Great Salt Lake composition that take $a_\mathrm{w}$ from about 0.98 to 0.75 at fixed ion
ratios, and swaps (Cl$^-$ for SO$_4^{2-}$, Na$^+$ for Mg$^{2+}$) designed in PHREEQC [4] to hold
$a_\mathrm{w}$ constant. Vials equilibrate with each sorbent at room temperature and on a
heater-shaker at 40--65\,$^\circ$C. ICP-OES, on gravimetric, matrix-matched dilutions, measures Li,
Na, K, Mg and B, and any Ti or Al the sorbent sheds. Equilibrium is confirmed by time series and by
approach from both sides (loading and stripping). Bayesian experimental design picks each next batch
where H1 and H2 predict the most different $K$: AI used to tell mechanisms apart, not to optimize
uptake. (Automated recovery from produced water so far targets selective precipitation [5].)

\wpflow
  {\textbf{Design}\\ next brines where H1 and H2 predict the most different $K$}
  {\textbf{OT-2}\\ brines by mass + sorbent; sealed vials, RT--65\,$^\circ$C}
  {\textbf{ICP-OES}\\ Li, Na, K, Mg, B; time series; load and strip}
  {\textbf{Pitzer/PHREEQC}\\ activities $\rightarrow$ exchange constants; fit H1, H2}
  {\textbf{Outcome}\\ waters released per exchange; role of the co-ion}
  {update the probability of H1 and H2, then choose the next batch}

\hd{Expected outcomes.} (1) Activity-based Li$^+$/Na$^+$ and Li$^+$/Mg$^{2+}$ exchange constants,
with enthalpies from their temperature dependence, for two sorbents with different uptake chemistry;
(2) the number of waters released per exchange, and whether co-ion pairing adds to it; (3) a test of
whether those constants predict selectivity in real Great Salt Lake and produced-water samples. The
result is a mechanistic basis for predicting selectivity in brines not yet measured, a step toward
domestic lithium from brines.

\hd{Fit and resources.} This addresses \emph{Overcoming Thermodynamics via Cooperative Chemical
Interactions}: a dilute critical mineral isolated from a complex mixture, with co-ion cooperation one
of the two mechanisms tested and AI used for mechanistic insight. No sorbent screening, optimization
or process engineering is proposed. The PI works on self-driving labs and Bayesian experimental
design; the lab's OT-2 already runs precipitation series checked against PHREEQC/Pitzer; campus
ICP-OES and XRD are available.

\begin{wprefs}
[1] C.\,E. Harvie, N. M{\o}ller, J.\,H. Weare, \emph{Geochim.\ Cosmochim.\ Acta} 48, 723 (1984),
doi:10.1016/0016-7037(84)90098-X.
[2] Y. Marcus, \emph{J.\ Chem.\ Soc.\ Faraday Trans.} 87, 2995 (1991), doi:10.1039/FT9918702995.
[3] M.\,P. Paranthaman \emph{et al.}, \emph{Environ.\ Sci.\ Technol.} 51, 13481 (2017),
doi:10.1021/acs.est.7b03464.
[4] D.\,L. Parkhurst, C.\,A.\,J. Appelo, USGS Techniques and Methods 6-A43 (2013), doi:10.3133/tm6A43.
[5] Ritchhart \emph{et al.}, \emph{Mater.\ Horiz.} 13, 7008 (2026), doi:10.1039/d6mh00475j.
\end{wprefs}

\end{document}
